\documentclass[11pt,a4paper]{article}
\usepackage[UTF8,fontset=none]{ctex}
\usepackage[margin=2.5cm]{geometry}
\usepackage[numbers,sort&compress]{natbib}
\usepackage{xcolor}
\usepackage{hyperref}
\hypersetup{colorlinks=true,linkcolor=blue,citecolor=blue,urlcolor=cyan}
\title{中文预训练语料质量与训练效用：机制、方法与迁移边界}
\author{}
\date{\today}

\begin{document}
\maketitle

\begin{abstract}
大语言模型性能不仅取决于参数和训练 token 数量，也取决于语料中信息、噪声、重复、领域覆盖与语言多样性的组合。现有研究已从启发式清洗推进到模型打分、数据价值估计、动态重加权和混合优化，但“质量分数是否能够预测训练效用”仍缺少跨来源、跨模型的稳定证据。高质量筛选通常能提高单位 token 收益，却可能同步删除长尾知识、特殊语体和低资源语言；同一选择方法在不同训练阶段也可能产生不同作用。因而，下一步研究应把质量表征、学习动力学和能力变化置于统一的受控实验框架，分别识别质量增益、覆盖损失与分布迁移，并报告效应大小、不确定性和失效条件。对中文语料而言，还需检验英文网页语料上形成的质量尺度能否迁移，以及文字规范、分词特征和语体差异怎样改变评价结果。
\end{abstract}

\section{引言}

语言模型训练长期遵循“模型、数据和计算共同扩展”的路径。早期规模律将损失与模型规模、数据规模和计算量联系起来，为训练预算配置提供了可量化坐标\cite{scalinglawsfor2020}；计算最优分析进一步指出，在固定计算量下，模型参数和训练 token 应协同增长\cite{trainingcomputeoptimal2022}。然而，这类规律通常把 token 当作相对同质的计量单位，较少刻画来源、语义密度、重复度和语言结构差异。Gopher 等大模型训练实践已经将质量过滤和数据混合作为关键环节\cite{scalinglanguagemodels2021}，由此引出一个更基础的问题：同样数量的 token 为什么产生不同训练收益？

开放语料的发展使该问题能够在更透明的条件下研究。RefinedWeb 表明，经系统清洗和去重的网页语料可以形成有竞争力的训练基础\cite{therefinedwebdataset2023}；Dolma 公开了多来源语料、处理链和工具，支持复现实验\cite{dolmaanopen2024}；BaichuanSEED 则强调广泛采集与去重对模型基线的影响\cite{baichuanseedsharingthe2024}。这些工作证明了数据工程的重要性，却没有自动回答“哪些质量特征通过何种机制影响哪些能力”。要把经验配方转成科学知识，需要区分数据的可用性、表面质量、训练贡献和外部能力效用，并用等 token、等算力和跨分布评价建立因果上更清晰的证据链。

\section{语料质量的概念分解与可测量性}

预训练语料质量不是单一属性。数据管理综述通常把采集、清洗、去重、过滤、混合与治理视为连续链条\cite{datamanagementfor2023}，数据集综述则显示不同模型公开的来源和处理细节差异很大\cite{datasetsforlarge2024}。因此，一个质量分数可能同时混合可读性、知识性、来源可信度、毒性、语言识别置信度和领域偏好。若不说明构成与标注目标，分数升高并不能说明训练效用必然改善。

语料的档案属性同样重要：抓取日期、网页生命周期、授权变化和删除请求决定了数据是否可追溯、可复核\cite{anarchivalperspective2024}。这类前提条件与学术意义上的“效用”应分开处理。含有隐私或不可合法使用的样本，即使对损失下降有贡献，也不能因效用高而进入训练集。生产模型中的训练数据抽取研究表明，记忆和泄漏风险并非纯理论问题\cite{extractingtrainingdata2020}；在更大规模模型上，数据抽取仍可发生\cite{scalableextractionof2023}。因此，合规、隐私与来源可追溯性应作为准入约束，而不是与效用分数加权后相互抵消的普通维度。

对训练相关质量，可把指标分成三层。第一层是文档内在属性，如结构完整、语言连贯、重复片段和信息密度；第二层是集合属性，如主题覆盖、语体多样性与样本间冗余；第三层是模型条件属性，即样本相对于当前模型、训练阶段和目标能力的边际贡献。MiniPile 等数据高效实验说明，小规模代表性子集能够用于快速比较数据策略\cite{theminipilechallenge2023}，但代理模型上的排序能否迁移到更大模型仍需单独验证。大规模语料审计也提示，只有知道数据中实际包含什么，质量讨论才有可靠对象\cite{whatsin2023}。

\section{质量筛选、去重与覆盖损失}

去重是较稳定的数据处理环节。对语言模型训练集进行文档与序列级去重能够减少记忆，并改善部分评价表现\cite{deduplicatingtrainingdata2022}。但去重阈值不是无害参数：过松会保留模板化网页和近重复内容，过严则可能删除合法复述、跨来源共识和低频表达。因而需要同时报告压缩率、重复簇结构、语言与领域保留率，而不能只报告训练集缩小比例。

质量过滤更依赖评价器的隐含偏好。基于百科、书籍或高教育价值页面训练的分类器，容易把正式书面语作为“好数据”的代理；这可能改善标准基准，却降低口语、方言和边缘领域覆盖。有关有害语言数据的系统分析早已指出，训练数据定义和标注规范会塑造模型偏差\cite{directionsinabusive2020}；从预训练数据追踪政治偏差的研究进一步说明，上游分布差异能够沿模型传播到下游任务\cite{frompretrainingdata2023}。所以，过滤策略必须以多维保留率和能力谱为约束，而非追求平均质量分数最大化。

多样性导向的数据选择试图缓解这一矛盾。相关方法在高质量候选中加入覆盖或代表性目标，以避免把相似样本反复选入\cite{harnessingdiversityfor2024}。网页重写研究则显示，改变低质量网页的表达而不是简单丢弃，也可能提高计算与数据效率\cite{rephrasingtheweb2024}。两者共同提示，数据治理至少包含“删除、保留、重写、降权”四种操作；何种操作最优取决于缺陷是否可修复、样本是否承载稀缺信息以及处理成本。

合成与增强数据使边界更复杂。大模型既可充当标注器和生成器，也可能放大原模型偏差或降低分布真实性\cite{largelanguagemodels2024}。合成数据策展综述强调生成、筛选和评价应形成闭环\cite{onllmsdriven2024}；数据增强综述则表明，不同生成范式的收益高度依赖任务和质量控制\cite{dataaugmentationusing2024}。因此，合成样本不宜与自然语料混合后只报告总体结果，应分别估计其增益、偏差和覆盖变化。

\section{数据价值与学习动力学机制}

如果内在质量描述“样本看起来怎样”，数据价值则试图回答“样本对模型产生什么作用”。训练数据影响分析综述将相关方法概括为重训练、影响函数、梯度相似性、表示近邻和数据归因等路线\cite{trainingdatainfluence2024}。OpenDataVal 提供了统一比较不同数据价值算法的基准\cite{opendatavalaunified2023}，而无需完整训练模型的数据价值方法试图降低估计成本\cite{datavaluationwithout2023}。这些进展为机制研究提供工具，但估计值通常依赖目标验证集；一旦验证集不能代表真实能力目标，数据价值就会退化成对局部指标的优化。

影响的时间结构尤其关键。研究发现，语言数据在训练早期与后期的影响估计并不等价\cite{firstisbetter2022}。自影响重加权据此用样本对自身损失或梯度的作用调整预训练权重\cite{selfinfluenceguided2023}；动态损失重加权则进一步随训练状态改变样本优先级\cite{dynamiclossbased2025}。Skill-it 将能力看作由训练数据中的技能及其依赖关系共同形成，并用数据驱动方式安排训练\cite{skillita2023}。这些方法共同表明，固定质量分数忽略了模型“已经学会什么”和“下一阶段缺什么”。

记忆、学习难度和效用也不能简单等同。跨模型记忆研究显示，模型规模、数据重复和提示条件会影响可提取记忆\cite{quantifyingmemorizationacross2022}。ORCA 等方法尝试定位提示输出所依赖的预训练证据\cite{orcainterpretingprompted2022}，机器遗忘也被用于反向检验特定预训练数据的作用\cite{decipheringtheimpact2024}。更大规模的事实追踪方法正在把影响估计扩展到预训练场景\cite{scalableinfluenceand2024}，LLM 规模影响函数则探索在现实参数规模下估计数据价值\cite{whatisyour2025}。然而，单样本归因仍受非凸优化、参数交互和近重复证据影响，适合生成机制假说，不宜未经干预实验便当作因果结论。

\section{数据混合、规模规律与训练效用}

预训练集通常由多个领域组成，混合比例决定模型在不同能力之间的资源分配。DoReMi 用小型代理模型估计领域损失并优化混合权重，在相同训练预算下提高效率\cite{doremioptimizingdata2023}。DoGE 直接用泛化估计进行领域重加权，强调目标分布而非单一训练损失\cite{dogedomainreweighting2023}。数据混合律则试图预测不同配比下的语言建模表现，减少昂贵的穷举训练\cite{datamixinglaws2024}；双层优化方法把验证目标嵌入数据重加权过程\cite{scalebioscalablebilevel2024}。这些方法将“语料好坏”转成条件化的资源配置问题。

但代理模型迁移是共同风险。较小模型能够廉价扫描大量配比，却可能与大模型在能力涌现、记忆容量和领域迁移上不同。SAVA 等学习无关数据价值方法追求更好的规模扩展性\cite{savascalablelearning2024}，多样性选择也试图降低对特定模型的依赖\cite{harnessingdiversityfor2024}。评估这类方法时，应至少交叉改变模型规模、训练预算与数据来源，并比较排序相关、绝对效应和最优阈值漂移。只在同一语料、同一模型家族上报告提升，不能证明评价方法可迁移。

规模律在此提供基线而非答案。若过滤后 token 数减少，性能变化可能同时来自质量升高、有效多样性下降和训练轮次增加。公平比较应明确采用等 token、等独立样本还是等计算量。数据选择综述指出，不同研究的预算口径和下游任务差异阻碍横向比较\cite{asurveyon2024}；数据中心 AI 立场研究也强调，模型时代的数据问题需要独立的方法学框架\cite{positionpaperdata2024}。因此，训练效用应报告为多目标向量：单位 token 收益、单位计算收益、能力覆盖、稳健性、污染风险和数据处理成本，而不是单一准确率。

\section{中文语料与跨来源、跨模型迁移}

中文语料评价不能直接照搬英文网页分类器。汉字与分词方式、简繁转换、古今语体、网络混写和中英夹杂都会改变困惑度、长度和可读性指标。ChatGLM 系列展示了中文与双语模型在数据、训练和工具能力上的持续扩展\cite{chatglmafamily2024}，BaichuanSEED 的数据构建经验也说明中文及多语网页需要独立的采集、过滤与去重设计\cite{baichuanseedsharingthe2024}。这些技术报告可提供工程基线，但公开细节仍不足以独立识别各质量环节的因果贡献。

评价端同样存在偏差。CMMLU 提供中文多任务理解测量\cite{cmmlumeasuringmassive2024}，但任何固定基准都可能受到训练污染、题型偏好和反复调参影响。对大模型基准适切性的批评指出，单一排行榜越来越难覆盖真实能力与可靠性\cite{inadequaciesoflarge2024}。因此，中文训练效用评价应同时包含留出语言建模集、人工构造的语言规范与语体测试、知识和推理任务、长尾与领域覆盖，以及污染审计；外部测试集不能反复参与质量阈值和混合权重选择。

跨来源验证应把“同指标不同含义”作为默认假设。例如，高困惑度在破损网页中可能表示噪声，在方言、少数语体或专业文本中却可能表示稀缺信息。跨模型验证则要区分共同趋势与尺度敏感效应：一个适合小模型的高频基础语料配比，未必适合容量更大的模型。Falcon 系列和 RefinedWeb 的经验说明开放网页能够支撑有竞争力的模型\cite{thefalconseries2023}，但这种结论不能自动外推到中文语料或不同 tokenizer。可迁移评价的目标不是找到永恒阈值，而是给出需要重新校准的条件、校准样本量和不确定性。

\section{讨论}

现有证据最清楚的共识是：数据处理会显著改变训练结果，去重、来源配比和质量过滤都不是可忽略的实现细节。主要分歧在于，研究者常用不同语料、模型、tokenizer、预算和基准，导致“质量提升”的含义不一致。部分工作测量训练损失，部分工作测量少数下游任务，还有工作以人工偏好或来源标签作为质量代理；这些指标之间相关但不等价。

机制研究面临三类识别问题。第一，质量维度共变：高可读文档往往也更长、更正式、来源更集中。第二，筛选改变训练分布和重复率，无法把结果简单归因于被优化的那一个分数。第三，模型状态形成反馈：模型打分器的偏好被用于筛选后，又可能在新模型中被强化。因而，可靠实验需要采用分层抽样、正交或近正交干预、等预算对照、多随机种子和预先登记的评价协议，并同时报告失败结果。

中文场景还存在资源与有效性的双重约束。大规模从零预训练虽接近真实应用，却会使实验数量受到计算预算限制；只用极小模型则可能失去外推意义。可行路径是先用代理模型和小规模干预筛除明显无效假说，再选择少量关键结论在更大模型上复核。若项目无法证明全国范围或所有领域的普适性，就应把成果表述为带边界的评价方法，而不是统一标准。

\section{展望}

未来值得优先推进四类可检验研究。其一，建立同一批中文文档的多维人工标注与自动指标，测量一致性、重测可靠性和跨来源漂移。其二，对若干质量维度做受控扰动，结合梯度统计、学习顺序和遗忘行为，识别从数据特征到能力变化的中介路径。其三，在等 token 与等算力两种口径下比较规则、模型打分、数据价值和混合优化方法，绘制质量收益与覆盖损失的帕累托前沿。其四，跨语料来源、模型尺度和训练阶段验证，给出校准程序、置信区间和拒绝外推条件。

开放科学也是方法的一部分。公开数据索引、处理配置、随机种子、代码版本和失败实验，能减少数据配方不可见造成的复现困难。数据生命周期研究强调，数据应在采集、处理、训练、评价和退出环节保持可追踪\cite{adatacentric2025}。对许可受限的数据，可发布索引、统计与处理脚本，而不重新分发原始文本。只有把证据链和边界同时公开，质量评价才可能从一次性配方发展为可累积研究。

\section{结论}

预训练语料质量研究正在从静态清洗走向条件化效用建模。去重和基础过滤提供必要底座，数据价值、学习动力学与混合优化则揭示样本贡献随模型、阶段和目标变化。当前最大的知识缺口不是再发明一个综合分数，而是解释质量特征如何影响学习过程、何时因覆盖损失而失效，以及这种关系能否跨来源和模型迁移。面向中文语料，合理的研究框架应把多维质量表征、受控训练、能力谱评价和迁移校准连成闭环，并以不确定性和适用边界约束结论。

\bibliographystyle{gbt7714-nsfc}
\bibliography{llm_corpus_quality_参考文献}
\end{document}
